\documentclass[10pt]{article}
\usepackage[T1]{fontenc}
\usepackage[utf8]{inputenc}
\usepackage[french]{babel}
\usepackage{amsmath,amssymb,mathrsfs,mathpazo}
\usepackage{geometry}
\geometry{a4paper,margin=2cm}
\usepackage{graphicx,array,tabularx,tabularray,colortbl,multirow,multicol,booktabs,enumitem,pifont,ulem,fancybox,diagbox,cancel}
\usepackage{pstricks,pstricks-add,pst-plot,pst-tree,pst-node,pst-text,pst-func,pst-math,pst-bspline,pst-3dplot,pst-eucl,pst-solides3d}
\let\pscalculate\undefined
\usepackage{pst-bezier}
\usepackage{tikz,pgfplots,tkz-tab}
\usetikzlibrary{arrows,arrows.meta,patterns,calc,positioning,trees,decorations.markings,shapes.geometric,shapes.misc,shapes.arrows,intersections,decorations.pathmorphing,decorations.pathreplacing,backgrounds,shadows,babel}
\pgfplotsset{compat=1.18}
\usepgfplotslibrary{fillbetween}
\usepackage[np]{numprint}
\usepackage{listings,makecell,pas-tableur,twoopt}
\let\barre\undefined
\usepackage[pstricks]{tablvar}
\definecolor{darkgreen}{RGB}{0,100,0}
\providecommand{\red}{\color{red}}
\providecommand{\blue}{\color{blue}}
\providecommand{\green}{\color{green}}
\providecommand{\black}{\color{black}}
\providecommand{\white}{\color{white}}
\providecommand{\gray}{\color{gray}}
\providecommand{\cyan}{\color{cyan}}
\providecommand{\magenta}{\color{magenta}}
\providecommand{\yellow}{\color{yellow}}
\providecommand{\violet}{\color{violet}}
\providecommand{\orange}{\color{orange}}
\providecommand{\euro}{\texteuro}
\providecommand{\EUR}{\texteuro}
\providecommand{\ieme}{\textsuperscript{e}}
\providecommand{\degres}{\ensuremath{^\circ}}
\providecommand{\degree}{\ensuremath{^\circ}}
\providecommand{\vv}[1]{\overrightarrow{#1}}
\providecommand{\Coord}[1]{#1}
\providecommand{\decosix}{}
\providecommand{\starredbullet}{\ensuremath{\bullet}}
\pagestyle{empty}
\setlength{\parindent}{0pt}
\usepackage{adjustbox}\DeclareUnicodeCharacter{2236}{:}\DeclareUnicodeCharacter{2190}{\ensuremath{\leftarrow}}
\providecommand{\R}{\mathbb{R}}
\providecommand{\N}{\mathbb{N}}
\providecommand{\Z}{\mathbb{Z}}
\providecommand{\Q}{\mathbb{Q}}
\providecommand{\D}{\mathbb{D}}
\providecommand{\C}{\mathbb{C}}
\providecommand{\e}{{\mathrm{e}}}
\providecommand{\vect}[1]{\overrightarrow{#1}}
\providecommand{\vectt}[1]{\overrightarrow{\mathrm{#1}}}
\providecommand{\vvt}[1]{\overrightarrow{\mathrm{#1}}}
\providecommand{\np}[2][]{{#2~#1}}
\providecommand{\vv}[1]{{\overrightarrow{#1}}}
\renewcommand{\d}{{\mathrm{d}}}
\providecommand{\degres}{{^\circ}}
\providecommand{\degre}{{^\circ}}
\providecommand{\blue}{\color{blue}}
\providecommand{\red}{\color{red}}
\providecommand{\enskip}{\quad}
\providecommand{\llbracket}{\lbrack\!\lbrack}
\providecommand{\rrbracket}{\rbrack\!\rbrack}
\providecommand{\euro}{{\text{€}}}
\providecommand{\no}{{\text{n°}}}
\providecommand{\interval}[2]{{[#1;#2]}}
\providecommand{\intervFO}[2]{{[#1;#2[}}
\providecommand{\dotfill}{\ldots}

\providecommand{\ds}{\displaystyle}
\providecommand{\textding}[1]{#1}
\providecommand{\Oij}{\(\mathrm{O};\vec{\imath},\vec{\jmath}\)}
\providecommand{\Oijk}{\(\mathrm{O};\vec{\imath},\vec{\jmath},\vec{k}\)}
\providecommand{\Ouv}{\(\mathrm{O};\vec{u},\vec{v}\)}
\providecommand{\intervFF}[2]{\ensuremath{[#1;#2]}}
\renewcommand{\euro}{\ifmmode\text{\texteuro}\else\texteuro\fi}
\DeclareUnicodeCharacter{00B0}{\ifmmode^{\circ}\else\textdegree\fi}
\DeclareUnicodeCharacter{20AC}{\ifmmode\text{\texteuro}\else\texteuro\fi}
% Énoncés : Éducation nationale. Transcription et corrigés : APMEP, Denis Vergès et contributeurs. Découpage et classement : CoopMaths, Sébastien Lozano, Éric Elter et contributeurs. Adaptation de la présentation : Maths Explorer.
% CC BY-SA : https://coopmaths.fr/www/about/
% https://coopmaths.fr/alea/static/bac/2026/tex/bac_2026_06_sujet1_polynesie_1_cor.tex
\begin{document}
\noindent\textit{APMEP / CoopMaths. CC BY-SA. Présentation adaptée par Maths Explorer.}\par\medskip
\medskip

Une école de surf travaille avec des planches en mousse pour les cours débutants, et d'autres en résine pour les cours de perfectionnement. Les planches en mousse représentent 80~\% du stock de l'école.

\begin{list}{\textbullet}{La gérante observe qu'après chaque saison,}
\item 75\,\% des planches en mousse sont en bon état et réutilisables la saison suivante. Les autres sont recyclées (et donc non réutilisables)~;
\item 60\,\% des planches en résine sont en bon état, les autres sont soit consolidées, soit recyclées.
\end{list}

Après une saison, on choisit une planche au hasard et on s'intéresse aux évènements suivants~:

\begin{tabular}{@{~~\textbullet~~} l @{~~~~\textbullet~~} l}
$M$ : \og La planche est en mousse \fg~;&
$B$ : \og La planche est en bon état \fg~;\\
$R$ : \og La planche sera recyclée \fg~;&
$C$ : \og La planche sera consolidée \fg.\\
\end{tabular}

\subsection*{Partie A}

\begin{enumerate}
\item On complète l'arbre pondéré ci-dessous avec les données de l'énoncé.



\begin{center}





















\pstree[levelsep=3.5,treesep=.7,treemode=R,nodesepB=3pt,nodesepA=0pt,nrot=:U]{\TR{}}
{
	\pstree[nodesepA=3pt]{\TR{$M$}\naput{$0,8$}}
	{
	\TR{$B$}\naput{$0,75$}
	\TR{$R$}\nbput{$1-0,75=0,25$}
	}
	\pstree[nodesepA=3pt]{\TR{$\overline{M}$}\nbput{$1-0,8=0,2$}}
	{
	\TR{$B$}\naput{$0,6$}
	\TR{$R$}
	\TR{$C$}
	}
}

\end{center}

\item On sait que $P\left(\overline{M} \cap R\right) = 0,05$.

\begin{enumerate}[label=\textbf{\alph*.}]
\item
Cela signifie que la probabilité que la planche choisie au hasard soit en résine et à recycler est égale à $0,05$.

\item On a choisi une planche en résine. 
La probabilité que cette planche soit consolidée pour la saison suivante est: $P_{\overline{M}}\left (C\right )$.

On sait que $P\left(\overline{M} \cap R\right) = 0,05$.
Mais $P\left(\overline{M} \cap R\right) = P\left (\overline{M}\right )\times P_{\overline{M}}\left (R\right )$, 

donc $P_{\overline{M}}\left (R\right ) = \dfrac{P\left(\overline{M} \cap R\right)}{P\left (\overline{M}\right ).} = \dfrac{0,05}{0,2}= 0,25$.

D'après les propriétés des probabilités, on a:

$P_{\overline{M}} \left (B\right ) +P_{\overline{M}} \left (R\right )  +P_{\overline{M}} \left (C\right )  =1$ donc 
$0,6 + 0,25 + P_{\overline{M}} \left (C\right )  =1$

et donc $P_{\overline{M}} \left (C\right ) =1-0,6-0,25=0,15$.
\end{enumerate}

 On peut donc compléter l'arbre pondéré.
 
 \begin{center}
 \pstree[levelsep=3.5,treesep=.7,treemode=R,nodesepB=3pt,nodesepA=0pt,nrot=:U]{\TR{}}
{
	\pstree[nodesepA=3pt]{\TR{$M$}\naput{$0,8$}}
	{
	\TR{$B$}\naput{$0,75$}
	\TR{$R$}\nbput{$1-0,75=0,25$}
	}
	\pstree[nodesepA=3pt]{\TR{$\overline{M}$}\nbput{$1-0,8=0,2$}}
	{
	\TR{$B$}\naput{$0,6$}
	\TR{$R$}\ncput*{\blue $0,25$}
	\TR{$C$}\ncput*{\blue $0,15$}
	}
}
\end{center}

\item La gérante affirme qu'après une saison, au moins $\dfrac{3}{4}$ des planches sont en bon état.


Il faut donc chercher $P\left (B\right )$.
D'après la formule des probabilités totales:

$P\left (B\right ) = P\left (M \cap B\right ) + P\left (\overline{M} \cap B\right ) = 0,8\times 0,75 + 0,2 \times 0,6=0,72<0,75$

Donc la gérante a tort.

\end{enumerate}

\bigskip

\subsection*{Partie B}

La gérante du club de surf envisage de proposer à certains clients une remise pour la saison suivante.

Pour cela, elle imagine le jeu suivant qui sera proposé à chaque client~:

Une urne contient 5 boules noires et un nombre $n$ de boules rouges.
Les boules sont indiscernables au toucher.
Le client tire simultanément deux boules dans l'urne.
Si les deux boules sont de même couleur, le client gagne une remise, sinon, il perd le jeu et n'obtient pas de remise.

\begin{enumerate}
\item 


Le client obtient une remise dans deux cas disjoints: il tire deux boules noires, ou il tire deux boules rouges (si $n\geqslant 2$).

\begin{list}{\textbullet}{}
\item Il y a $n+5$ boules en tout: le nombre de façons de tirer simultanément 2 boules parmi $n+5$ est $\binom{n+5}{2}$.

\item Il y a $5$ boules noires: le nombre de façons de tirer simultanément 2 boules noires parmi $5$ est $\binom{5}{2}$.

Les boules sont indiscernables au toucher donc il y a équiprobabilité; la probabilité de tirer 2 boules noires est donc $p_1=\dfrac{\binom{5}{2}}{\binom{n+5}{2}}$.

\item Il y a $n$ boules rouges: le nombre de façons de tirer simultanément 2 boules rouges parmi $n$ est $\binom{n}{2}$.

Les boules sont indiscernables au toucher donc il y a équiprobabilité; la probabilité de tirer 2 boules rouges est donc $p_2=\dfrac{\binom{n}{2}}{\binom{n+5}{2}}$.
\end{list}

Les évènements \og tirer 2 boules noires \fg{} et \og tirer 2 boules rouges\fg{}  sont disjoints donc la probabilité de tirer 2 boules de même couleur est:

$p=p_1+p_2= \dfrac{\binom{5}{2}}{\binom{n+5}{2}} + \dfrac{\binom{n}{2}}{\binom{n+5}{2}} = \dfrac{\binom{5}{2} +\binom{n}{2}}{\binom{n+5}{2}}$.

\item 


On sait que $\ds\binom{n}{p} = \dfrac{n\,!}{p\,! \times (n-p)\,!}$.

On en déduit que:
$\ds\binom{5}{2}=\dfrac{5\,!}{3\,! \times 2\,!}= \dfrac{5\times 4}{2}=10$,
que: $\ds\binom{n}{2}=\dfrac{n\,!}{\left (n-2\right )\,! \times 2\,!}= \dfrac{n\left (n-1\right )}{2}$,\\
et que: $\ds\binom{n+5}{2}=\dfrac{(n+5)\,!}{\left (n+3\right )\,! \times 2\,!}= \dfrac{\left (n+5\right )\left (n+4\right )}{2}$.

Donc
$p= \dfrac{10 + \frac{n\left (n-1\right )}{2}}{\frac{\left (n+5\right )\left (n+4\right )}{2}}
= \dfrac{\frac{20}{2} + \frac{n\left (n-1\right )}{2}}{\frac{\left (n+5\right )\left (n+4\right )}{2}}
=\dfrac{20+n\left (n-1\right )}{\left (n+5\right )\left (n+4\right )}$.

\item La gérante souhaite s'assurer que la probabilité de gagner est inférieure à $\dfrac{1}{2}$.



On cherche donc $n$ tel que $p<\dfrac{1}{2}$; on résout cette inéquation.

$\begin{aligned}
p<\dfrac{1}{2}
& \iff \dfrac{20+n\left (n-1\right )}{\left (n+5\right )\left (n+4\right )}<\dfrac{1}{2}
 \iff \dfrac{20+n\left (n-1\right )}{\left (n+5\right )\left (n+4\right )} - \dfrac{1} {2} <0\\
 & \iff \dfrac{2\left (20+n\left (n-1\right )\right ) - \left (n+5\right )\left (n+4\right )}{2\left (n+5\right )\left (n+4\right )}<0
\iff \dfrac{40+2n^2-2n-n^2-5n-4n-20}{2\left (n+5\right )\left (n+4\right )} < 0\\
& \iff \dfrac{n^2-11n +20}{2\left (n+5\right )\left (n+4\right )} < 0
 \iff n^2-11n+20<0 \text{~~car~~} 2\left (n+5\right )\left (n+4\right )>0
\end{aligned}$

On étudie le signe du trinôme $x^2-11x+20$.

$\Delta = (-11)^2-4\times 1 \times 20=121-80=41>0$;
donc ce trinôme admet deux racines:

$x'=\dfrac{-(-11) -\sqrt{41}}{2\times 1}= \dfrac{11-\sqrt{41}}{2}\approx 2,3$ et $x''= \dfrac{11+\sqrt{41}}{2}\approx 8,7$ 

On en déduit le signe du trinôme.

\begin{center}
$ \begin{tablvar}[intervalwidth=7em,stretch=1.2]{3}
 \hline
 x & -\infty &  & x' & & x'' &  & +\infty\\
 \hline
ax^2+bx+c & & \text{signe de } a & \barre[0] & \text{signe de } -a  & \barre[0] & \text{signe de } a  & \\
 \hline
x^2-11x+20  & & + & \barre[0] & -  & \barre[0] & + & \\
 \hline
 \end{tablvar}$

\end{center}
L'ensemble des valeurs possibles du nombre $n$ de boules rouges qu'elle peut mettre dans l'urne pour que la probabilité de remise est donc l'ensemble des nombres entiers entre $x'\approx 2,3$ et $x''\approx 8,7$; c'est donc l'ensemble $\left \{ 3\::\: 4 \::\: 5 \::\: 6 \::\: 7 \::\: 8\strut \right \}$.

\end{enumerate}

\bigskip

\end{document}
